%! Setting Up a Test for a Population Mean — Unit 7, Lesson 7.4 — Homework
\documentclass[10pt]{article}
\usepackage{apstats-article}
\usepackage{apstats-boxes}
\usepackage{grading}
\gradingassignment{apstats-u07-l04-hw}
\gradingcourse{advanced_statistics}
\newcommand{\TallMath}[1]{$\displaystyle #1\rule[-1.4em]{0pt}{3.2em}$}

\begin{document}

\pageheader{Unit 7, Lesson 7.4}{Homework}
{\small\itshape Write each final answer in its box. Only what is inside a box is graded.}\par

\vspace{0.15in}

\textit{For each problem: (a) name the inference method and df; (b) state $H_0$ and $H_a$;
(c) verify both conditions. Show all work.}

\begin{enumerate}

\item \textbf{Water Temperature.} \quad
A marine biologist believes the mean water temperature at a reef has risen above the historical
average of 77$^\circ$F. She randomly collects 28 temperature readings; the data are slightly
right-skewed with no outliers. The sample standard deviation is 2.3$^\circ$F.

\begin{enumerate}[label=\alph*.]
  \item Method: \makebox[5.5cm]{\smash{\gradedblank[id=hw1a1, type=short, points=1, answer={{one-sample t-test for mu}}, check=text, display={\boldmath one-sample $t$-test for $\mu$}]{5.3cm}{0.19in}}} \quad $df = $ \makebox[1.5cm]{\smash{\gradedblank[id=hw1a2, type=short, points=1, answer=27, check=numeric, tolerance=0]{1.3cm}{0.19in}}}
  \item $H_0$: \makebox[4.5cm]{\smash{\gradedblank[id=hw1b1, type=short, points=1, answer={{mu = 77}}, check=text, display={\boldmath$\mu = 77$}]{4.3cm}{0.19in}}} \quad $H_a$: \makebox[4.5cm]{\smash{\gradedblank[id=hw1b2, type=short, points=1, answer={{mu > 77}}, check=text, display={\boldmath$\mu > 77$}]{4.3cm}{0.19in}}}
  \item Conditions:
        \begin{itemize}
          \boxguard
          \item Random: \hfill\makebox[4.6in][l]{\gradedhere{\gradedblank[id=hw1c1, type=rubric, points=1, display={28 readings randomly collected --- met}]{4.6in}{0.19in}}}\\
                \gradedhide{\writelines{9}}
          \boxguard
          \item Normal ($n < 30$): \hfill\makebox[4.6in][l]{\gradedhere{\gradedblank[id=hw1c2, type=rubric, points=1, display={$n = 28 < 30$, slight skew, no outliers --- met}]{4.6in}{0.19in}}}\\
                \gradedhide{\writelines{9}}
        \end{itemize}
\end{enumerate}
\vspace{0.1in}

\newpage
\item \textbf{Delivery Times.} \quad
An online retailer claims the mean delivery time for orders is 3 days. A consumer group
believes the true mean is \emph{different} from 3 days. They randomly audit 55 recent orders.

\begin{enumerate}[label=\alph*.]
  \item Method: \makebox[5.5cm]{\smash{\gradedblank[id=hw2a1, type=short, points=1, answer={{one-sample t-test for mu}}, check=text, display={\boldmath one-sample $t$-test for $\mu$}]{5.3cm}{0.19in}}} \quad $df = $ \makebox[1.5cm]{\smash{\gradedblank[id=hw2a2, type=short, points=1, answer=54, check=numeric, tolerance=0]{1.3cm}{0.19in}}}
  \item $H_0$: \makebox[4.5cm]{\smash{\gradedblank[id=hw2b1, type=short, points=1, answer={{mu = 3}}, check=text, display={\boldmath$\mu = 3$}]{4.3cm}{0.19in}}} \quad $H_a$: \makebox[4.5cm]{\smash{\gradedblank[id=hw2b2, type=short, points=1, answer={{mu != 3}}, check=text, display={\boldmath$\mu \ne 3$}]{4.3cm}{0.19in}}}
  \item Conditions:
        \begin{itemize}
          \boxguard
          \item Random: \hfill\makebox[4.6in][l]{\gradedhere{\gradedblank[id=hw2c1, type=rubric, points=1, display={random audit of 55 orders --- met}]{4.6in}{0.19in}}}\\
                \gradedhide{\writelines{9}}
          \boxguard
          \item Normal: \hfill\makebox[4.6in][l]{\gradedhere{\gradedblank[id=hw2c2, type=rubric, points=1, display={$n = 55 \ge 30$, CLT --- met}]{4.6in}{0.19in}}}\\
                \gradedhide{\writelines{9}}
        \end{itemize}
\end{enumerate}
\vspace{0.1in}

\newpage
\item \textbf{Matched Pairs: Reaction Time.} \quad
A researcher tests whether a new app reduces reaction time. Ten participants each complete
a reaction-time task \emph{before} and \emph{after} using the app for two weeks. The
differences (before $-$ after) show a roughly symmetric distribution with no outliers.

\begin{enumerate}[label=\alph*.]
  \item Define the order of subtraction and the parameter $\mu_d$:

        $d_i = $ \makebox[3cm]{\smash{\gradedblank[id=hw3a1, type=short, points=1, answer=before, check=text]{2.8cm}{0.19in}}} $-$ \makebox[3cm]{\smash{\gradedblank[id=hw3a2, type=short, points=1, answer=after, check=text]{2.8cm}{0.19in}}}

        $\mu_d = $\par
        \gradedhide{\gradedhere{\gradedblank[id=hw3a3, type=rubric, points=1, display={true mean difference (before $-$ after)}]{\linewidth+6pt}{0.19in}}\writelines{9}}

  \item Method: \makebox[5cm]{\smash{\gradedblank[id=hw3b1, type=short, points=1, answer={{matched-pairs t-test for mu_d}}, check=text, display={matched-pairs $t$-test for $\mu_d$}]{4.8cm}{0.19in}}} \quad $df = $ \makebox[1.5cm]{\smash{\gradedblank[id=hw3b2, type=short, points=1, answer=9, check=numeric, tolerance=0]{1.3cm}{0.19in}}}
  \item $H_0$: \makebox[4cm]{\smash{\gradedblank[id=hw3c1, type=short, points=1, answer={{mu_d = 0}}, check=text, display={\boldmath$\mu_d = 0$}]{3.8cm}{0.19in}}} \quad $H_a$: \makebox[4cm]{\smash{\gradedblank[id=hw3c2, type=short, points=1, answer={{mu_d > 0}}, check=text, display={\boldmath$\mu_d > 0$}]{3.8cm}{0.19in}}}
  \item Conditions:
        \begin{itemize}
          \boxguard
          \item Random/Independence: \hfill\makebox[4.6in][l]{\gradedhere{\gradedblank[id=hw3d1, type=rubric, points=1, display={random assignment / representative sample; participants independent}]{4.6in}{0.19in}}}\\
                \gradedhide{\writelines{9}}
          \boxguard
          \item Normal ($n < 30$): \hfill\makebox[4.6in][l]{\gradedhere{\gradedblank[id=hw3d2, type=rubric, points=1, display={$n = 10 < 30$, differences symmetric, no outliers --- met}]{4.6in}{0.19in}}}\\
                \gradedhide{\writelines{9}}
        \end{itemize}
\end{enumerate}
\vspace{0.1in}

\newpage
\item \textbf{AP-Style.} \quad
\textit{A researcher collects a random sample of 12 observations to test $H_0: \mu = 50$ vs.\
$H_a: \mu \ne 50$. She verifies that the dotplot of the data shows moderate right skewness and
no outliers. Which of the following best describes the Normal condition?}

\begin{enumerate}[label=(\Alph*)]
  \item The condition is met because $n = 12$ is large enough for the CLT.
  \item The condition is not met; moderate right skewness with $n < 30$ violates the condition.
  \item The condition is met; moderate skewness with no outliers is acceptable for $n < 30$.
  \item The condition cannot be evaluated without knowing the population standard deviation.
\end{enumerate}
\vspace{0.06in}
Answer: \makebox[1.5cm]{\smash{\gradedblank[id=hw4a, type=choice, points=1, answer=C]{0.4in}{0.2in}}}\rule[-0.22in]{0pt}{0.52in}

Brief justification:\par
\gradedhide{\gradedhere{\gradedblank[id=hw4b, type=rubric, points=1, display={$n < 30$: moderate skew with no outliers is acceptable}]{\linewidth+6pt}{0.19in}}\writelines{9}}

\end{enumerate}

\begin{extensionbox}
A statistician says: ``When $n$ is small and the data are not Normal, I can use a sign test or
Wilcoxon signed-rank test instead of a $t$-test.'' Research what a \textbf{nonparametric test}
is and explain in 2--3 sentences why these alternatives exist and what trade-off they involve.
\writelines{3}
\end{extensionbox}

\end{document}
